\chapter{Raízes quadradas}\label{ch:g9:sqrt}

Qual é o comprimento da diagonal de um quadrado de lado $1$? O teorema
de Pitágoras responde $\sqrt 2$, um número cujo quadrado é $2$ --- e
que, no fim das contas, não é uma \omterm{def:g9:fractions:fraction}{fração}. Este capítulo define as
\omterm{def:g9:sqrt:def}{raízes quadradas}, estabelece as regras de cálculo com elas e ensina a
simplificar expressões como $\sqrt{75}$.

\section{Definição e primeiras propriedades}

\begin{definition}[Raiz quadrada]\label{def:g9:sqrt:def}
Seja $a \geq 0$. A \emph{raiz quadrada}\index{raiz quadrada} de $a$,
escrita $\sqrt a$, é o único número \emph{não negativo} cujo quadrado é
$a$:
\[
\sqrt a \geq 0
\qquad\text{e}\qquad
\left(\sqrt a\right)^2 = a .
\]
Números negativos não têm raiz quadrada, já que todo quadrado é não
negativo.
\end{definition}

\begin{example}
$\sqrt{49} = 7$, $\sqrt 0 = 0$, $\sqrt 1 = 1$, $\sqrt{2.25} = 1.5$. As
primeiras \omterm{def:g9:sqrt:def}{raízes quadradas} a saber de cor são as dos \emph{quadrados
perfeitos}: $1, 4, 9, 16, 25, 36, 49, 64, 81, 100, 121, 144$.
\end{example}

\begin{omfigure}
\begin{tikzpicture}[scale=1.35]
  \draw[thick, fill=omDef!8] (0,0) rectangle (2,2);
  \draw[omThm, thick] (0,0) -- (2,2);
  \node[below, font=\small] at (1,0) {$1$};
  \node[left, font=\small] at (0,1) {$1$};
  \node[omThm, above left=-2pt, font=\small, rotate=45] at (1.08,0.92)
    {$\sqrt2$};
  \draw[thin] (1.75,0) -- (1.75,0.25) -- (2,0.25);
\end{tikzpicture}

{\small A diagonal de um quadrado de lado $1$ tem comprimento
$\sqrt{1^2 + 1^2} = \sqrt2 \approx 1.414$, pelo teorema de Pitágoras (o
\cref{ch:g9:trig} o relembra). Esse número é irracional: os decimais
dele nunca se repetem.}
\end{omfigure}

\begin{proposition}[Raiz quadrada de um quadrado]\label{prop:g9:sqrt:square}
Para todo número real $x$ (positivo ou não):
\[
\sqrt{x^2} = \abs{x} =
\begin{cases}
x & \text{se } x \geq 0,\\
-x & \text{se } x < 0 .
\end{cases}
\]
\end{proposition}

\begin{proof}
O número $\abs{x}$ é não negativo e o quadrado dele é $x^2$ (um número
e o oposto dele têm o mesmo quadrado). Sendo o único número não
negativo de quadrado $x^2$, ele é $\sqrt{x^2}$.
\end{proof}

\begin{example}
$\sqrt{(-5)^2} = \sqrt{25} = 5 = \abs{-5}$: a \omterm{def:g9:sqrt:def}{raiz quadrada}
``esquece'' o sinal, não o restaura. Em particular, $\sqrt{x^2} = x$ é
\emph{errado} para $x$ negativo.
\end{example}

\section{Produtos e quocientes de raízes quadradas}

\begin{theorem}[Regras de multiplicação e divisão]\label{thm:g9:sqrt:rules}
Para todos os $a \geq 0$ e $b \geq 0$:
\[
\sqrt{a b} = \sqrt a \times \sqrt b,
\qquad\text{e, para } b > 0:\quad
\sqrt{\frac ab} = \frac{\sqrt a}{\sqrt b} .
\]
\end{theorem}

\begin{proof}
O número $\sqrt a \times \sqrt b$ é não negativo (\omterm{def:g2:mult:def}{produto} de números
não negativos), e o quadrado dele é
\[
\left(\sqrt a \times \sqrt b\right)^2
= \left(\sqrt a\right)^2 \times \left(\sqrt b\right)^2
= a b .
\]
Pela unicidade do número não negativo de quadrado $ab$, ele é igual a
$\sqrt{ab}$. A regra do \omterm{def:g3:division:remainder}{quociente} se demonstra do mesmo jeito.
\end{proof}

\begin{remark}[Não existe regra dessas para somas!]
$\sqrt{a + b}$ \emph{não} é $\sqrt a + \sqrt b$ em geral:
\[
\sqrt{9 + 16} = \sqrt{25} = 5,
\qquad\text{mas}\qquad
\sqrt 9 + \sqrt{16} = 3 + 4 = 7 .
\]
\end{remark}

\begin{method}[Simplificar $\sqrt n$]\label{met:g9:sqrt:simplify}
Para simplificar a \omterm{def:g9:sqrt:def}{raiz quadrada} de um inteiro:
\begin{enumerate}
  \item encontre o maior quadrado perfeito que \omterm{def:g9:arith:divisor}{divide} $n$ (decomponha
        $n$, se precisar);
  \item separe com a regra do \omterm{def:g2:mult:def}{produto} e extraia o quadrado:
        $\sqrt{k^2 m} = k\sqrt m$;
  \item confira que nenhum quadrado perfeito \omterm{def:g9:arith:divisor}{divide} o que sobrou sob a
        raiz.
\end{enumerate}
\end{method}

\begin{example}\label{ex:g9:sqrt:simplify}
Simplifique $\sqrt{75}$: como $75 = 25 \times 3$,
\[
\sqrt{75} = \sqrt{25 \times 3} = \sqrt{25} \times \sqrt 3 = 5\sqrt3 .
\]
Simplifique $\sqrt{72}$: como $72 = 36 \times 2$,
$\sqrt{72} = 6\sqrt2$. E um \omterm{def:g3:division:remainder}{quociente}:
$\sqrt{\dfrac{49}{16}} = \dfrac{\sqrt{49}}{\sqrt{16}} = \dfrac74$.
\end{example}

\begin{example}[Somar raízes quadradas]\label{ex:g9:sqrt:add}
\omterm{def:g1:addition:def}{Somas} de \omterm{def:g9:sqrt:def}{raízes quadradas} só simplificam quando as raízes são
\emph{parecidas}. Calcule $\sqrt{12} + \sqrt{27} - \sqrt{48}$,
simplificando cada parcela primeiro:
\begin{align*}
\sqrt{12} &= \sqrt{4 \times 3} = 2\sqrt3, \\
\sqrt{27} &= \sqrt{9 \times 3} = 3\sqrt3, \\
\sqrt{48} &= \sqrt{16 \times 3} = 4\sqrt3,
\end{align*}
então a \omterm{def:g1:addition:def}{soma} é
$2\sqrt3 + 3\sqrt3 - 4\sqrt3 = (2 + 3 - 4)\sqrt3 = \sqrt3$.
\end{example}

\begin{example}[Desenvolver com raízes quadradas]\label{ex:g9:sqrt:expand}
As identidades da álgebra se aplicam às \omterm{def:g9:sqrt:def}{raízes quadradas}. Desenvolva
$\left(\sqrt5 + 2\right)^2$ com $(a+b)^2 = a^2 + 2ab + b^2$:
\[
\left(\sqrt5 + 2\right)^2 = 5 + 2 \times 2\sqrt5 + 4 = 9 + 4\sqrt5 .
\]
E com o terceiro \omterm{def:g2:mult:def}{produto} notável, $(a+b)(a-b) = a^2 - b^2$:
\[
\left(\sqrt7 + \sqrt3\right)\left(\sqrt7 - \sqrt3\right) = 7 - 3 = 4 :
\]
o \omterm{def:g2:mult:def}{produto} de dois números irracionais pode ser um inteiro.
\end{example}

\section{A equação \texorpdfstring{$x^2 = a$}{x2 = a}}

\begin{theorem}[Resolver $x^2 = a$]\label{thm:g9:sqrt:equation}
\begin{enumerate}
  \item Se $a > 0$, a \omterm{def:g8:equations:def}{equação} $x^2 = a$ tem exatamente duas soluções:
        $\sqrt a$ e $-\sqrt a$;
  \item se $a = 0$, a única solução é $0$;
  \item se $a < 0$, não há solução.
\end{enumerate}
\end{theorem}

\begin{proof}
Reescreva $x^2 = a$ como $x^2 - \left(\sqrt a\right)^2 = 0$ quando
$a \geq 0$ e fatore como \omterm{ex:g1:subtraction:difference}{diferença} de quadrados:
\[
\left(x - \sqrt a\right)\left(x + \sqrt a\right) = 0 ,
\]
então $x = \sqrt a$ ou $x = -\sqrt a$ (as duas coincidem quando
$a = 0$). Quando $a < 0$ não há solução, já que $x^2 \geq 0 > a$ para
todo $x$.
\end{proof}

\begin{omfigure}
\begin{tikzpicture}
\begin{axis}[omaxis,
  width=7.6cm, height=5.4cm,
  xmin=-3.1, xmax=3.1, ymin=-1.6, ymax=5.4,
  xtick={-1.732,1.732}, xticklabels={$-\sqrt3$,$\sqrt3$},
  ytick={3}, yticklabels={$3$}]
  \addplot[omDef, thick, domain=-2.25:2.25, samples=90] {x^2};
  \addplot[omThm, thick, domain=-2.9:2.9] {3};
  \fill (axis cs:-1.732,3) circle (1.5pt);
  \fill (axis cs:1.732,3) circle (1.5pt);
  \draw[dashed] (axis cs:-1.732,3) -- (axis cs:-1.732,0);
  \draw[dashed] (axis cs:1.732,3) -- (axis cs:1.732,0);
\end{axis}
\end{tikzpicture}

{\small Resolver $x^2 = 3$ sobre a parábola $y = x^2$: a reta
horizontal $y = 3$ a corta nas duas abscissas $\pm\sqrt3$.}
\end{omfigure}

\begin{example}
$x^2 = 16$: soluções $4$ e $-4$. \quad
$x^2 = 5$: soluções $\sqrt5$ e $-\sqrt5$ (valores exatos;
$\sqrt5 \approx 2.236$). \quad
$x^2 = -9$: nenhuma solução. \quad
$3x^2 = 21$: divida por 3 primeiro, $x^2 = 7$, soluções $\pm\sqrt7$.
\end{example}

\section{Exercícios}

\begin{exercise}[$\star$]\label{exo:g9:sqrt:1}
Calcule sem calculadora:
\[
\sqrt{64}, \qquad \sqrt{100}, \qquad \sqrt{0.09}, \qquad
\sqrt{\tfrac{1}{25}}, \qquad \left(\sqrt{13}\right)^2, \qquad
\sqrt{(-4)^2} .
\]
\end{exercise}

\begin{exercise}[$\star$]\label{exo:g9:sqrt:2}
Simplifique:
\[
\sqrt{50}, \qquad \sqrt{45}, \qquad \sqrt{98}, \qquad \sqrt{300} .
\]
\end{exercise}

\begin{exercise}[$\star$]\label{exo:g9:sqrt:3}
Calcule e simplifique:
\[
\sqrt2 \times \sqrt{18}, \qquad
\frac{\sqrt{75}}{\sqrt3}, \qquad
\sqrt5 \times \sqrt{20}, \qquad
\frac{\sqrt{8}}{\sqrt{50}} .
\]
\end{exercise}

\begin{exercise}[$\star$]\label{exo:g9:sqrt:4}
Resolva: $x^2 = 36$; \quad $x^2 = 11$; \quad $x^2 + 4 = 0$; \quad
$2x^2 = 50$.
\end{exercise}

\begin{exercise}[$\star$]\label{exo:g9:sqrt:5}
Reduza a um único termo: $\sqrt{20} + \sqrt{45}$ e depois
$3\sqrt8 - \sqrt{18} + \sqrt2$.
\end{exercise}

\begin{exercise}[$\star\star$]\label{exo:g9:sqrt:6}
Desenvolva e simplifique:
\[
\left(\sqrt3 + 1\right)^2, \qquad
\left(2\sqrt5 - 3\right)^2, \qquad
\left(\sqrt6 + \sqrt2\right)\left(\sqrt6 - \sqrt2\right).
\]
\end{exercise}

\begin{exercise}[$\star\star$]\label{exo:g9:sqrt:7}
Um campo quadrado tem \omterm{def:g6:measure:area}{área} $6400$ m$^2$. Qual é o comprimento do lado
dele? E o da diagonal (valor exato e depois arredondado ao metro)?
\end{exercise}

\begin{exercise}[$\star\star$]\label{exo:g9:sqrt:8}
Mostre que $\dfrac{1}{\sqrt2} = \dfrac{\sqrt2}{2}$ (multiplique
\omterm{def:g4:fractions:def}{numerador} e \omterm{def:g4:fractions:def}{denominador} por $\sqrt2$). Use o mesmo truque para escrever
$\dfrac{6}{\sqrt3}$ e $\dfrac{10}{\sqrt5}$ sem \omterm{def:g9:sqrt:def}{raiz quadrada} no
\omterm{def:g4:fractions:def}{denominador}.
\end{exercise}

\begin{exercise}[$\star\star$]\label{exo:g9:sqrt:9}
Verdadeiro ou falso? Justifique com uma demonstração ou um
contraexemplo.
\begin{enumerate}
  \item Para todos os $a, b \geq 0$: $\sqrt{ab} = \sqrt a \sqrt b$.
  \item Para todos os $a, b \geq 0$: $\sqrt{a+b} = \sqrt a + \sqrt b$.
  \item Para todo $x$: $\sqrt{x^2} = x$.
  \item $\left(3\sqrt2\right)^2 = 18$.
\end{enumerate}
\end{exercise}

\begin{exercise}[$\star\star\star$]\label{exo:g9:sqrt:10}
Seja $x = \sqrt{7 + 4\sqrt3}$. Calcule $\left(2 + \sqrt3\right)^2$ e
deduza uma expressão mais simples de $x$. Mesma pergunta para
$\sqrt{9 - 4\sqrt5}$ (mire num quadrado da forma
$\left(\sqrt5 - 2\right)^2$ e cuidado com o sinal).
\end{exercise}

\section{Problema: o número que não é uma fração}

\begin{problem}[{Problema de fim de semana --- a demonstração lendária
de que $\sqrt2$ é irracional, os muitos parentes dele e a receita
espantosamente boa de Heron para aproximá-lo}]
\label{pb:g9:sqrt:1}
A diagonal de um quadrado de lado $1$ é um comprimento perfeitamente
real --- você a traçou no \cref{ex:g8:pythagoras:diagonal} --- e, mesmo
assim, como os pitagóricos descobriram com horror há uns vinte e cinco
séculos, \emph{\omterm{def:g9:fractions:fraction}{fração} nenhuma} a mede. A lenda conta que a descoberta
foi punida com afogamento. Este problema o conduz pela demonstração
imortal (quatro questões e ela é sua para a vida toda), multiplica as
vítimas e termina com a receita que os engenheiros usaram durante dois
mil anos para domar o número indomável.

\medskip
\noindent\textbf{Parte I --- A demonstração.}
\begin{enumerate}
  \item Primeiro cerque-o: verifique que $1.4 < \sqrt2 < 1.5$ e depois
        que $1.41 < \sqrt2 < 1.42$, elevando os limites ao quadrado.
        Mais um algarismo: entre que números de três casas decimais
        $\sqrt2$ se encontra?
  \item Agora suponha --- para chegar a uma contradição --- que
        $\sqrt2 = \frac ab$ para alguma \omterm{def:g9:fractions:fraction}{fração} \emph{irredutível}
        (\cref{met:g9:arith:simplify}). Eleve os dois lados ao quadrado
        e mostre que $a^2 = 2 b^2$. Qual é a paridade de $a^2$?
  \item Os fatos de paridade do \cref{pb:g8:pythagoras:1} dizem: os
        \omterm{def:g3:numbers:evenodd}{números ímpares} têm quadrados \omterm{def:g3:numbers:evenodd}{ímpares}. Deduza que $a$ é par,
        escreva $a = 2k$, substitua --- e mostre que $b$ também tem de
        ser par.
  \item Onde está a contradição? Conclua e enuncie o teorema por
        inteiro: \emph{$\sqrt2$ é irracional --- não é \omterm{def:g3:division:remainder}{quociente} de
        números inteiros}.
  \item A demonstração se apoiou uma vez, discretamente, na palavra
        ``irredutível''. Aponte o passo exato que desabaria sem ela e
        explique por que supor a \omterm{def:g9:fractions:fraction}{fração} irredutível era legítimo, para
        começo de conversa.
\end{enumerate}

\medskip
\noindent\textbf{Parte II --- As vítimas se multiplicam.}
\begin{enumerate}[resume]
  \item Um segundo estilo de demonstração, pelas decomposições em
        fatores primos (\cref{thm:g9:arith:factorization}): em
        $a^2 = 3b^2$, compare a paridade do \emph{\omterm{def:g8:powers:def}{expoente} do $3$} de
        cada lado (\cref{pb:g9:arith:1}: os quadrados carregam
        \omterm{def:g8:powers:def}{expoentes} pares). Conclua que $\sqrt3$ é irracional.
  \item Rode o mesmo argumento dos \omterm{def:g8:powers:def}{expoentes} em $a^2 = n b^2$ para um
        número inteiro $n$ qualquer: para que $n$ ele produz uma
        contradição, e para quais ele falha? Enuncie o resultado
        completo: $\sqrt n$ é irracional exatamente quando\,\dots
  \item Demonstre o lema do escudo: um racional não nulo vezes um
        irracional é irracional. (Suponha $r x = q$ com $r, q$
        racionais, $r \neq 0$, e isole $x$.) Deduza que $2\sqrt2$ e
        $\frac{\sqrt2}{2}$ são irracionais.
  \item Demonstre que $1 + \sqrt2$ é irracional. E depois a bonita:
        supondo que $s = \sqrt2 + \sqrt3$ fosse racional, calcule
        $s^2$, isole $\sqrt6$ e encontre a contradição.
  \item Modere o entusiasmo: dê dois números irracionais cuja
        \emph{\omterm{def:g1:addition:def}{soma}} é racional e dois cujo \emph{\omterm{def:g2:mult:def}{produto}} é racional.
        (A irracionalidade não se conserva pelas operações --- cada
        caso exige a própria demonstração.)
\end{enumerate}

\medskip
\noindent\textbf{Parte III --- A receita de Heron.}
Dois mil anos antes das calculadoras, Heron de Alexandria aproximava
$\sqrt2$ assim: \emph{chute $x$; troque o chute pela média de $x$ e
$\frac2x$; repita.}
\begin{enumerate}[resume]
  \item Parta do chute $x = 1$ e calcule os dois chutes seguintes como
        frações exatas.
  \item Calcule o terceiro chute, de novo como \omterm{def:g9:fractions:fraction}{fração} exata, e o valor
        decimal dele. Compare com a questão 1: quantas casas decimais
        de $\sqrt2$ já estão corretas?
  \item A ideia por trás da receita: um retângulo de \omterm{def:g6:measure:area}{área} $2$ com um
        lado $x$ tem o outro lado $\frac2x$. Mostre que $\sqrt2$ fica
        sempre \emph{entre} $x$ e $\frac2x$ (considere os dois casos
        $x^2 > 2$ e $x^2 < 2$), então tirar a média dos dois lados
        aperta o retângulo em direção ao quadrado.
  \item Os chutes $\frac32$, $\frac{17}{12}$, $\frac{577}{408}$
        escondem uma joia: calcule $3^2 - 2 \times 2^2$, depois
        $17^2 - 2 \times 12^2$, depois $577^2 - 2 \times 408^2$. A
        Parte I demonstrou que $a^2 - 2b^2 = 0$ é impossível --- quão
        perto do impossível chegam as frações de Heron?
  \item Final, em duas frases: a diagonal do quadrado de lado $1$
        existe no papel, \omterm{def:g9:fractions:fraction}{fração} nenhuma a mede, e a escrita decimal
        dela nunca pode se repetir (\cref{pb:g9:fractions:1}). Que tipo
        de ``números novos'' a reta numérica tem, portanto, de conter
        --- e em que ponto desta série a história completa deles é
        contada?
\end{enumerate}
\end{problem}
